\documentclass[11pt]{article}

\usepackage{fontspec}
\usepackage{xeCJK}
\usepackage{amsmath}
\usepackage{amssymb}
\usepackage{array}
\usepackage{booktabs}
\usepackage{longtable}
\usepackage[margin=1in]{geometry}
\usepackage[hidelinks]{hyperref}
\usepackage{seqsplit}

\setmainfont{Times New Roman}
\setCJKmainfont{Malgun Gothic}
\setCJKsansfont{Malgun Gothic}
\setCJKmonofont{Malgun Gothic}
\xeCJKsetup{CJKspace=true}
\XeTeXlinebreaklocale "ko"
\XeTeXlinebreakskip=0pt plus 1pt
\setlength{\emergencystretch}{2em}

\newcommand{\claimstatus}[2]{\textbf{#1 [#2]}}
\newcommand{\poly}{\operatorname{poly}}
\newcolumntype{L}[1]{>{\raggedright\arraybackslash}p{#1}}
\renewcommand{\abstractname}{초록}

\title{고전적 정수분해를 위한 차수 분리 지수:\\
세 개의 제한 정리와 열여덟 개의 구조적 장벽}
\author{MOSEF 연구 프로그램}
\date{2026년 7월 28일}

\begin{document}
\maketitle

\begin{abstract}
이 논문은 이진 입력 길이
\(m=\lceil\log_2N\rceil\)을 복잡도 기준으로 삼아, 고전적 정수분해를
위한 차수 분리 지수 계열을 연구한다. 일반 정수분해의 고전적 다항
시간 알고리즘이 존재한다고 가정하지 않는다. 저장소에서 검증된
결과는 인수분해에 의존하는 두 개의 유전적 약속 클래스에 대한
Las Vegas 정리, 하나의 완전한 유리 root-of-unity 분류 정리, 그리고
명시적 지수 목록, 압축 표현, 산술 회로, 기하급수 quotient, 예외적
cyclotomic cofactor 일정에 대한 구조적 장벽들이다.

이번 M29 결과는 길이별 공개 후보
\[
A=3,\qquad B_m=2^m+3,\qquad g=2
\]
의 \(\Phi_4\) cofactor \(C_m=C_4(2)\)를 정확히 분석한다. 다음 닫힌
식을 증명한다.
\[
C_m=\frac{16(2^{3\cdot2^m+5}+3)}{35}.
\]
소수 5와 7에는 분모 때문에 별도의 quotient 판정이 필요하고,
\(p>7\)에서는 \(p\mid C_m\)이 하나의 모듈러 합동식과 정확히
동치이다. 또한 \(\gcd(C_m,C_{m+1})=16\)이다. 한 cofactor가 \(s\)개의
소수 중 \(h\)개를 지지하면 정확히 \(h(s-h)\)개의 semiprime 쌍만
proper GCD를 만든다. 같은 쪽의 hit--hit 쌍은 full collision이고
miss--miss 쌍은 unit이다. 따라서 \(s\ge3\)인 모집단 전체를 한
후보가 분리할 수 없다. 이 결론은 여러 후보, 적응형 일정, 일반
산술 회로 또는 일반 정수분해에 대한 하한이 아니다.
\end{abstract}

\section{연구 목적과 주장 규율}

MOSEF는 ``Multi-Group Order-Separating Exponent Factoring'' 연구
방향의 이름이지, 이미 성립한 일반 알고리즘의 이름이 아니다. 모든
수학 문장은 다음 상태 중 하나로 관리한다.
\[
\texttt{DEFINITION},\ \texttt{PROVED},\ \texttt{CONDITIONAL},\
\texttt{CONJECTURE},\ \texttt{HEURISTIC},\ \texttt{EMPIRICAL},\
\texttt{OPEN},\ \texttt{REFUTED}.
\]
실험에서 성공한 사실은 정리로 승격하지 않으며, exact integer의
크기와 서로 다른 소인수의 수를 동일시하지 않는다. 미지의 인수
\(p,q\), \(p-1\), \(q-1\), 또는 인수에 의존하는 local order를
공개 생성기의 입력으로 몰래 사용하지 않는다.

영문 원고 \path{paper/main.tex}는 모든 과거 정리와 증명의 상세
기록이다. 이 한국어 동반 원고는 동일한 claim ID와 상태를 사용하고,
현재 M29 결과의 정의, 정리, 증명, 실험, 한계 및 재현 절차를
자기완결적으로 제공한다.

\section{모형과 기본 판정}

\subsection{입력 길이와 차수 분리}

\(P(N)\)을 \(N\)의 서로 다른 소인수 집합이라 하자. \(g\)가 \(N\)과
서로소이고 \(d>0\)일 때
\[
D_N(g,d)=\{p\in P(N):\operatorname{ord}_p(g)\mid d\}
\]
로 둔다. 이 집합이 공집합도 아니고 \(P(N)\) 전체도 아니면
\[
\gcd(g^d-1,N)
\]
은 nontrivial factor이다. 그러나 \(N\)이 prime power이면 support만
보는 정의로는 부족하다. 정확한 조건은 각 \(p^e\parallel N\)에
대하여
\[
a_p=\min(e,v_p(g^d-1))
\]
를 기록하고, 어떤 \(a_p\)는 양수이면서 모든 \(a_p\)가 \(e\)인 것은
아닌지 확인하는 것이다.

\subsection{복잡도 원장}

다항 시간은 \(N\)이 아니라 \(m=\lceil\log_2N\rceil\)에 대한 다항
시간을 뜻한다. 다음 비용을 모두 센다.

\begin{itemize}
\item 공개 base, exponent, parameter의 생성과 비트 길이;
\item 모듈러 곱셈, 거듭제곱, 역원 시도, GCD와 factor extraction;
\item product tree나 회로의 모든 명시적 출력;
\item exact lift 또는 prime-support certificate를 요구할 때의 실제
  출력 비트 수;
\item 재귀적으로 완전 분해할 때의 전체 노드와 성공 확률.
\end{itemize}

compact modular evaluator가 존재한다는 사실은 그 exact lift가 짧다는
뜻이 아니다. 반대로 exact lift가 길다는 사실은 서로 다른 소인수가
많다는 뜻도 아니다.

\section{M0--M28의 검증 결과 요약}

\subsection{제한적 양의 결과}

저장소에는 세 개의 핵심 양의 분류가 있다.

\begin{enumerate}
\item 인수 구조로 정의된 유전적 \(p-1\) semismooth-asymmetry 약속
  클래스에서는 fresh uniform base를 사용한 Las Vegas 완전 분해가
  expected polynomial bit complexity를 가진다.
\item 별도의 유전적 nonsplit \(p+1\) 약속 클래스에서는 정확한 Lucas
  root count와 fresh independent parameter를 사용한 두 번째 Las Vegas
  정리가 성립한다.
\item unequal depth-two signed numerator의 rational root-of-unity 비율은
  common-step 비율 \(-1\), \(\Phi_4\) 비율 \(1\), \(\Phi_6\) 비율
  \(2\)의 세 계열로 완전히 분류된다.
\end{enumerate}

앞의 두 약속 클래스 membership은 미지의 인수분해에 의존할 수
있으므로, 일반 입력에서 이를 인식하는 알고리즘을 증명한 것은
아니다.

\subsection{구조적 장벽}

검증된 장벽들은 다음 원칙을 반복해서 확인한다.

\begin{itemize}
\item divisor를 hit하는 것과 서로 다른 prime-factor signature를
  분리하는 것은 다르다.
\item \(P=a+a^{-1}\)인 자연스러운 Lucas map은 독립 channel이 아니라
  multiplicative residue의 valuation을 제곱하는 재표현이다.
\item 짧은 exponent 목록, multiplication-only straight-line program,
  same-base addition/subtraction program은 명시된 모형 밖의 정보를
  공짜로 만들지 못한다.
\item materialized product tree와 product-only DAG의 aggregate GCD는
  개별 성공을 가릴 수 있지만 없던 proper factor를 만들지 못한다.
\item geometric sum과 nested quotient의 compact 계산은 가능하지만,
  분모의 unit/proper/full 분기와 공개 multiplier 경로를 모두 세면
  새로운 무제한 factorization 경로가 되지 않는다.
\item signed cross-stage aggregation은 product-only 결론을 벗어나는
  실제 counterexample을 만들므로 별도의 산술 모형이 필요하다.
\item 예외적 \(\Phi_4,\Phi_6\) cofactor는 direct cyclotomic GCD가
  놓치는 proper factor를 실제로 추출하지만, 고정된 유한 공개
  schedule은 무한히 많은 square-free semiprime을 놓친다.
\end{itemize}

\subsection{M28: compact 원장과 materialized 원장의 분리}

길이 \(m\)에서 입력 \(N\)보다 먼저 공개 일정 \(\mathcal T_m\)을
고르는 모형을 생각한다. exact 값들을 실제로 출력하고 그 총 비트
수를 \(W_m\)이라 하면, balanced prime 모집단에서 \(h_m\)개의 서로
다른 소수를 지지하기 위해
\[
\left\lfloor\frac{m-1}{2}\right\rfloor h_m\le W_m
\]
가 필요하다. 그러나 현재 family는 \(O(m)\)-bit parameter와 compact
modular evaluation을 가지면서 exact cofactor 길이가 exponential이다.
따라서 이 materialized bound를 compact 실행 시간 하한으로 옮길 수
없다. M29는 exact 크기가 아니라 prime support와 GCD collision을
직접 분석한다.

\section{M29: compact cofactor prime-support 장벽}

\subsection{정의와 정확한 판정}

\textbf{DEF-029 [DEFINITION].}
\(m\ge2\)에서
\[
B_m=2^m+3,\qquad C_m=C_4(2)
\]
로 둔다. 유한 odd-prime 모집단 \(\mathcal P\)에 대하여
\[
H_m(\mathcal P)=\{p\in\mathcal P:p\mid C_m\}
\]
를 analytical support라 한다. 공개 알고리즘은 이 집합이나
\(C_m\)의 factorization을 입력받지 않는다. 기존 M26 공식을 사용해
\(C_m\bmod N\)만 계산한다.

\textbf{BAR-023 [PROVED].}
\[
E_m=3\cdot2^m+5
\]
라 하면
\[
C_m=\frac{16(2^{E_m}+3)}{35},\qquad
v_2(C_m)=4,\qquad 3\nmid C_m.
\]
또한
\[
5\mid C_m\Longleftrightarrow m\equiv2\pmod4,
\]
\[
7\mid C_m\Longleftrightarrow m\equiv2\pmod3,
\]
이고 모든 소수 \(p>7\)에 대하여
\[
p\mid C_m\Longleftrightarrow2^{E_m}\equiv-3\pmod p.
\]
연속한 두 cofactor는
\[
\gcd(C_m,C_{m+1})=16
\]
을 만족한다.

\(\mathcal P\)가 서로 다른 odd prime \(s\)개를 포함하고
\(h=|H_m(\mathcal P)|\)라 하자. \(pq\) 꼴의 모든 square-free 쌍
중에서 다음 개수가 정확하다.

\begin{center}
\begin{tabular}{L{0.28\linewidth}L{0.26\linewidth}L{0.32\linewidth}}
\toprule
두 소수의 signature & 쌍의 수 & GCD 결과 \\
\midrule
hit--miss & \(h(s-h)\) & proper factor \\
hit--hit & \(\binom h2\) & full collision \(pq\) \\
miss--miss & \(\binom{s-h}{2}\) & unit \(1\) \\
\bottomrule
\end{tabular}
\end{center}

따라서
\[
h(s-h)\le\left\lfloor\frac{s^2}{4}\right\rfloor
<\binom{s}{2}\qquad(s\ge3)
\]
이고, 한 compact cofactor는 모집단의 모든 pair를 factor할 수 없다.

\subsection{증명}

예외적 \(\Phi_4\) identity를 \(g=2\)에서 계산하면
\[
5C_m=S_3(2)+S_{B_m}(8)
=7+\frac{8^{B_m}-1}{7}.
\]
그러므로
\[
35C_m=8^{B_m}+48
=16(2^{3B_m-4}+3)
=16(2^{E_m}+3).
\]
\(m\ge2\)이면 \(E_m\equiv1\pmod4\)이고
\(E_m\equiv2\pmod3\)이다. 따라서 \(2^{E_m}+3\)은 5와 7로 나누어지고,
위 식은 정수이다. 이 quotient는 odd이므로 \(v_2(C_m)=4\)이다.
\(E_m\)은 odd이므로
\[
2^{E_m}+3\equiv-1+0\equiv2\pmod3,
\]
즉 \(3\nmid C_m\)이다.

분모의 5를 제거한 뒤에도 5가 남으려면
\[
2^{E_m}\equiv-3\pmod{25}
\]
이어야 한다. 2의 modulo 25 order는 20이고
\(2^{17}\equiv-3\pmod{25}\)이므로
\[
E_m\equiv17\pmod{20}.
\]
이는 \(2^m\equiv4\pmod{20}\), 다시
\(m\equiv2\pmod4\)와 동치이다.

7의 경우 \(8=1+7\)을 사용하면
\[
16(2^{E_m}+3)=8^{B_m}+48
\equiv1+7B_m+48
\equiv7B_m\pmod{49}.
\]
따라서 quotient 뒤에 7이 남는 조건은 \(7\mid B_m\), 즉
\(2^m\equiv4\pmod7\)이고, 이는 \(m\equiv2\pmod3\)이다. \(p>7\)에서는
16과 35가 모두 unit이므로 닫힌 식에서 generic congruence가 즉시
따라온다.

이제 consecutive GCD를 보자. 두 cofactor의 공통 2-adic 부분은
정확히 16이다. 3은 어느 쪽도 나누지 않고, 5와 7의 level 조건은
연속해서 성립하지 않는다. 어떤 \(p>7\)이 둘을 모두 나눈다고 하자.
\[
E_{m+1}=2E_m-5
\]
이므로 generic congruence에서
\[
-3\equiv2^{E_{m+1}}
=2^{2E_m-5}
\equiv\frac{(-3)^2}{32}
=\frac9{32}\pmod p.
\]
따라서 \(p\mid105\)여야 하는데 \(p>7\)과 모순이다. 그러므로
\(\gcd(C_m,C_{m+1})=16\)이다.

마지막으로 서로 다른 \(p,q\)에 대하여 \(\gcd(C_m,pq)\)가 proper인
경우는 정확히 둘 중 하나만 \(H_m\)에 속하는 경우이다. cut의 서로
다른 쪽을 고르는 방법은 \(h(s-h)\)개이다. 양쪽 모두 hit이면 GCD는
\(pq\), 양쪽 모두 miss이면 GCD는 1이다. 세 count와
\(\lfloor s^2/4\rfloor\) 상한이 모두 따라온다. 이 증명은 support를
공개 입력으로 사용하지 않으며, finite experiment의 zero-hit 관측도
사용하지 않는다.

\subsection{반박된 해석}

\textbf{REF-025 [REFUTED].}
exact cofactor가 exponential하게 길다는 사실이나 여러 level에서
관측한 소인수들을 합친 사실만으로 한 GCD의 universal extraction을
보장할 수 없다. 연속 level의 odd support는 서로 겹치지 않고, 한
level에서는 support가 넓어질수록 hit--hit full collision도 늘어난다.
여러 후보가 multi-bit signature를 만드는 경우는 M30의 별도
연구 문제이다.

\section{구현과 bit complexity}

\(B_m=2^m+3\)과 \(E_m=3\cdot2^m+5\)의 이진 표현 길이는 \(O(m)\)이다.
M26의 geometric-sum recurrence는 exact \(C_m\)을 만들지 않고
\(C_m\bmod N\)을 \(O(m)\)개의 모듈러 composition step으로 계산한다.
Python은 arbitrary-precision semantic oracle이고, Rust는 문서화된
\texttt{u64} 범위의 독립 구현이며, C\#은 \texttt{BigInteger}로 선택된
vector를 다시 계산한다.

소수 \(p\)가 주어진 local profile 함수는 증명과 differential test를
위한 분석 도구이다. 실제 factorization 후보는 \(p\)를 받지 않고
\[
\gcd(C_m\bmod N,N)
\]
만 계산한다. support enumeration, exact cofactor factorization,
unknown local roots는 실행 시간에 숨겨져 있지 않다.

\section{재현 가능한 실험}

\textbf{EMP-028 [EMPIRICAL].}
EXP-0028은 다음을 deterministic하게 검사했다.

\begin{itemize}
\item 23개 level에서 prime 20,000 이하의 52,026개 profile;
\item 51,934개 generic congruence와 92개 small-prime quotient case;
\item 49,742개 consecutive-support case;
\item 13개 exact closed form과 consecutive GCD;
\item input length 9부터 40까지 82,019개의 balanced prime;
\item input length 20까지 2,034개의 explicit pair outcome;
\item 17개 canonical vector에 대한 34개 Python/Rust/C\# 비교.
\end{itemize}

실패는 0개였다. 등록된 balanced 범위에서는 same-index support hit가
없었다. 특히 length 40의 22,394개 prime과 250,734,421개 pair는 이
한 후보에서 모두 unit으로 분류되었다. 이는 유한 범위의 empirical
결과일 뿐이며, 이후 모든 길이에서 support가 0이라는 정리가 아니다.

canonical summary SHA-256은
\[
\texttt{\seqsplit{8ca6c6310b64e56d37cfbc98caba9deddd02d5c35ea909870341aae8f23efb7a}}
\]
이다.

\section{한계와 다음 연구 문제}

BAR-023이 증명하지 않은 사항은 다음과 같다.

\begin{itemize}
\item balanced population에서 \(H_m\)의 asymptotic density;
\item \(H_m\)을 factorization 없이 명시적으로 출력하는 recognizer;
\item 여러 base 또는 여러 exceptional parameter tuple의 결합 효과;
\item \(N\)-dependent 또는 adaptive schedule의 하한;
\item repeated-prime 입력 전체에 대한 이번 single-cut count의 직접
  확장;
\item 일반 arithmetic circuit 또는 일반 고전적 정수분해의 하한.
\end{itemize}

다음 M30의 bounded question은 길이별 polynomial candidate list
\((C_{m,1},\ldots,C_{m,r})\)가 각 prime에 multi-bit signature를
부여할 때, 모든 prime pair를 분리하는 정확한 필요충분조건과 최소
candidate count를 구하는 것이다. 먼저 signature injectivity를
증명한 뒤에만 구체적인 parameter schedule을 제안한다.

\section{결론}

M29는 compact exact lift의 크기와 prime-factor separation을 분리했다.
현재 \(\Phi_4\) family는 polynomial-size descriptor와 compact evaluator를
가지지만, 그 한 후보는 결국 모집단을 hit와 miss의 두 부분으로
나누는 한 비트만 제공한다. proper factor는 cut을 가로지르는
semiprime에서만 나오고, 같은 쪽의 두 소수는 unit 또는 full
collision을 만든다. exact lift가 길다는 사실은 이 combinatorial
장벽을 제거하지 않는다.

일반 고전적 다항 시간 정수분해 알고리즘도, 그러한 알고리즘의
일반적 불가능성도 여기서 주장하지 않는다. 검증된 기여는 정확한
closed form, exceptional prime 판정, consecutive GCD, single-signature
pair count, 독립 구현, 재현 가능한 finite audit이다.

\section{재현 명령}

\begin{verbatim}
python scripts/validate_foundation.py
python scripts/check_publication.py
python -m unittest discover -s tests -v
python -m compileall -q python scripts tests
ruff check python tests scripts
mypy python
cargo fmt --all --check
cargo clippy --workspace --all-targets --all-features -- -D warnings
cargo test --workspace --all-features
dotnet build verification/csharp/MosefVerifier.csproj `
  --no-restore --configuration Release
python scripts/check_baseline_differential.py
python scripts/run_m29_compact_cofactor_prime_support_audit.py `
  --level-min 2 --level-max 24 --prime-max 20000 `
  --input-length-min 9 --input-length-max 40 `
  --pair-check-max 20 --exact-level-max 14
python scripts/check_m29_compact_cofactor_prime_support_differential.py
python scripts/generate_korean_claim_appendix.py
latexmk -xelatex -outdir=tmp/pdfs `
  -interaction=nonstopmode -halt-on-error paper/main.tex
latexmk -xelatex -jobname=main-ko -outdir=tmp/pdfs `
  -interaction=nonstopmode -halt-on-error paper/main-ko.tex
\end{verbatim}

\appendix
\section{전체 claim ID와 상태}

아래 표는 \path{research/PUBLICATION_CLAIMS.md}에서 생성된다. 상세한
statement, hypothesis, evidence, scope와 adversarial review 상태의
권위 있는 원장은 \path{research/CLAIMS.md}이다.

% Generated by scripts/generate_korean_claim_appendix.py.
% Do not edit this inventory by hand.
\begin{longtable}{L{0.34\linewidth}L{0.56\linewidth}}
\toprule
Claim 및 상태 & 권위 있는 근거 위치 \\
\midrule
\endfirsthead
\toprule
Claim 및 상태 & 권위 있는 근거 위치 \\
\midrule
\endhead
\claimstatus{DEF-001}{DEFINITION} & \texttt{CODEX.md section 4} \\
\claimstatus{DEF-002}{DEFINITION} & \texttt{research/proofs/M2-formal-specification.md} \\
\claimstatus{DEF-003}{DEFINITION} & \texttt{research/proofs/M2-formal-specification.md} \\
\claimstatus{DEF-004}{DEFINITION} & \texttt{research/proofs/THM-001-semismooth-promise.md} \\
\claimstatus{LEM-001}{PROVED} & \texttt{research/proofs/M2-formal-specification.md} \\
\claimstatus{LEM-002}{PROVED} & \texttt{research/proofs/M2-formal-specification.md} \\
\claimstatus{THM-001}{PROVED} & \texttt{research/proofs/THM-001-semismooth-promise.md} \\
\claimstatus{EXT-001}{PROVED} & \texttt{research/literature/BASELINE.md SRC-001} \\
\claimstatus{EXT-002}{PROVED} & \texttt{research/literature/BASELINE.md SRC-002} \\
\claimstatus{OPEN-001}{REFUTED} & \texttt{research/NEGATIVE\_RESULTS.md NR-001} \\
\claimstatus{OPEN-002}{OPEN} & \texttt{research/proofs/M2-formal-specification.md} \\
\claimstatus{OPEN-003}{OPEN} & \texttt{research/proofs/M2-formal-specification.md} \\
\claimstatus{REF-001}{REFUTED} & \texttt{research/NEGATIVE\_RESULTS.md NR-002} \\
\claimstatus{EMP-001}{EMPIRICAL} & \texttt{schemas/baseline-vectors-v1.json} \\
\claimstatus{EMP-002}{EMPIRICAL} & \texttt{research/experiments/EXP-0002-m2-separator-search.md} \\
\claimstatus{EMP-003}{EMPIRICAL} & \texttt{research/experiments/EXP-0003-m3-semismooth-search.md} \\
\claimstatus{DEF-005}{DEFINITION} & \texttt{research/proofs/BAR-001-divisor-cover-separation-gap.md} \\
\claimstatus{EXT-003}{CONDITIONAL} & \texttt{research/literature/SRC-004-umans-wang-divisor-conjecture.md} \\
\claimstatus{BAR-001}{PROVED} & \texttt{research/proofs/BAR-001-divisor-cover-separation-gap.md} \\
\claimstatus{REF-002}{REFUTED} & \texttt{research/NEGATIVE\_RESULTS.md NR-003} \\
\claimstatus{EMP-004}{EMPIRICAL} & \texttt{research/experiments/EXP-0004-m4-difference-cover-search.md} \\
\claimstatus{DEF-006}{DEFINITION} & \texttt{research/proofs/BAR-002-conjugate-channel-correlation.md} \\
\claimstatus{EXT-004}{PROVED} & \texttt{research/literature/SRC-005-williams-p-plus-one.md} \\
\claimstatus{BAR-002}{PROVED} & \texttt{research/proofs/BAR-002-conjugate-channel-correlation.md} \\
\claimstatus{REF-003}{REFUTED} & \texttt{research/NEGATIVE\_RESULTS.md NR-004} \\
\claimstatus{EMP-005}{EMPIRICAL} & \texttt{research/experiments/EXP-0005-m5-multigroup-correlation.md} \\
\claimstatus{DEF-007}{DEFINITION} & \texttt{research/proofs/THM-002-nonsplit-lucas-promise.md} \\
\claimstatus{LEM-003}{PROVED} & \texttt{research/proofs/THM-002-nonsplit-lucas-promise.md} \\
\claimstatus{THM-002}{PROVED} & \texttt{research/proofs/THM-002-nonsplit-lucas-promise.md} \\
\claimstatus{EMP-006}{EMPIRICAL} & \texttt{research/experiments/EXP-0006-m7-nonsplit-lucas.md} \\
\claimstatus{DEF-008}{DEFINITION} & \texttt{research/proofs/BAR-003-combined-promise-density.md} \\
\claimstatus{BAR-003}{PROVED} & \texttt{research/proofs/BAR-003-combined-promise-density.md} \\
\claimstatus{REF-004}{REFUTED} & \texttt{research/NEGATIVE\_RESULTS.md NR-005} \\
\claimstatus{EMP-007}{EMPIRICAL} & \texttt{research/experiments/EXP-0007-m8-promise-density.md} \\
\claimstatus{DEF-009}{DEFINITION} & \texttt{research/proofs/BAR-004-exponent-encoding-divisor-budget.md} \\
\claimstatus{BAR-004}{PROVED} & \texttt{research/proofs/BAR-004-exponent-encoding-divisor-budget.md} \\
\claimstatus{REF-005}{REFUTED} & \texttt{research/NEGATIVE\_RESULTS.md NR-006} \\
\claimstatus{EMP-008}{EMPIRICAL} & \texttt{research/experiments/EXP-0008-m9-divisor-budget.md} \\
\claimstatus{DEF-010}{DEFINITION} & \texttt{research/proofs/BAR-005-multiplication-straight-line-compression.md} \\
\claimstatus{BAR-005}{PROVED} & \texttt{research/proofs/BAR-005-multiplication-straight-line-compression.md} \\
\claimstatus{REF-006}{REFUTED} & \texttt{research/NEGATIVE\_RESULTS.md NR-007} \\
\claimstatus{EMP-009}{EMPIRICAL} & \texttt{research/experiments/EXP-0009-m10-compressed-exponents.md} \\
\claimstatus{EXT-005}{PROVED} & \texttt{research/literature/SRC-006-rosser-schoenfeld-primes.md} \\
\claimstatus{DEF-011}{DEFINITION} & \texttt{research/proofs/BAR-006-boundary-constant.md} \\
\claimstatus{BAR-006}{PROVED} & \texttt{research/proofs/BAR-006-boundary-constant.md} \\
\claimstatus{REF-007}{REFUTED} & \texttt{research/NEGATIVE\_RESULTS.md NR-008} \\
\claimstatus{EMP-010}{EMPIRICAL} & \texttt{research/experiments/EXP-0010-m11-boundary-schedule.md} \\
\claimstatus{EXT-006}{PROVED} & \texttt{research/literature/SRC-007-lichtman-shifted-smooth-primes.md} \\
\claimstatus{DEF-012}{DEFINITION} & \texttt{research/proofs/BAR-007-primorial-factor-scale.md} \\
\claimstatus{BAR-007}{PROVED} & \texttt{research/proofs/BAR-007-primorial-factor-scale.md} \\
\claimstatus{REF-008}{REFUTED} & \texttt{research/NEGATIVE\_RESULTS.md NR-009} \\
\claimstatus{EMP-011}{EMPIRICAL} & \texttt{research/experiments/EXP-0011-m12-primorial-scale.md} \\
\claimstatus{DEF-013}{DEFINITION} & \texttt{research/proofs/BAR-008-general-factor-scale.md} \\
\claimstatus{BAR-008}{PROVED} & \texttt{research/proofs/BAR-008-general-factor-scale.md} \\
\claimstatus{REF-009}{REFUTED} & \texttt{research/NEGATIVE\_RESULTS.md NR-010} \\
\claimstatus{EMP-012}{EMPIRICAL} & \texttt{research/experiments/EXP-0012-m13-general-factor-scale.md} \\
\claimstatus{DEF-014}{DEFINITION} & \texttt{research/proofs/BAR-009-addition-subtraction.md} \\
\claimstatus{BAR-009}{PROVED} & \texttt{research/proofs/BAR-009-addition-subtraction.md} \\
\claimstatus{REF-010}{REFUTED} & \texttt{research/NEGATIVE\_RESULTS.md NR-011} \\
\claimstatus{EMP-013}{EMPIRICAL} & \texttt{research/experiments/EXP-0013-m14-addition-subtraction.md} \\
\claimstatus{DEF-015}{DEFINITION} & \texttt{research/proofs/BAR-010-implicit-batch.md} \\
\claimstatus{BAR-010}{PROVED} & \texttt{research/proofs/BAR-010-implicit-batch.md} \\
\claimstatus{REF-011}{REFUTED} & \texttt{research/NEGATIVE\_RESULTS.md NR-012} \\
\claimstatus{EMP-014}{EMPIRICAL} & \texttt{research/experiments/EXP-0014-m15-implicit-batch.md} \\
\claimstatus{DEF-016}{DEFINITION} & \texttt{research/proofs/BAR-011-product-dag.md} \\
\claimstatus{BAR-011}{PROVED} & \texttt{research/proofs/BAR-011-product-dag.md} \\
\claimstatus{REF-012}{REFUTED} & \texttt{research/NEGATIVE\_RESULTS.md NR-013} \\
\claimstatus{EMP-015}{EMPIRICAL} & \texttt{research/experiments/EXP-0015-m16-product-dag.md} \\
\claimstatus{DEF-017}{DEFINITION} & \texttt{research/proofs/BAR-012-dyadic-telescope.md} \\
\claimstatus{BAR-012}{PROVED} & \texttt{research/proofs/BAR-012-dyadic-telescope.md} \\
\claimstatus{REF-013}{REFUTED} & \texttt{research/NEGATIVE\_RESULTS.md NR-014} \\
\claimstatus{EMP-016}{EMPIRICAL} & \texttt{research/experiments/EXP-0016-m17-dyadic-telescope.md} \\
\claimstatus{DEF-018}{DEFINITION} & \texttt{research/proofs/BAR-013-arbitrary-geometric-sum.md} \\
\claimstatus{BAR-013}{PROVED} & \texttt{research/proofs/BAR-013-arbitrary-geometric-sum.md} \\
\claimstatus{REF-014}{REFUTED} & \texttt{research/NEGATIVE\_RESULTS.md NR-015} \\
\claimstatus{EMP-017}{EMPIRICAL} & \texttt{research/experiments/EXP-0017-m18-geometric-sum.md} \\
\claimstatus{DEF-019}{DEFINITION} & \texttt{research/proofs/BAR-014-nested-geometric-quotient.md} \\
\claimstatus{BAR-014}{PROVED} & \texttt{research/proofs/BAR-014-nested-geometric-quotient.md} \\
\claimstatus{REF-015}{REFUTED} & \texttt{research/NEGATIVE\_RESULTS.md NR-016} \\
\claimstatus{EMP-018}{EMPIRICAL} & \texttt{research/experiments/EXP-0018-m19-nested-quotient.md} \\
\claimstatus{DEF-020}{DEFINITION} & \texttt{research/proofs/BAR-015-iterated-geometric-quotient.md} \\
\claimstatus{BAR-015}{PROVED} & \texttt{research/proofs/BAR-015-iterated-geometric-quotient.md} \\
\claimstatus{REF-016}{REFUTED} & \texttt{research/NEGATIVE\_RESULTS.md NR-017} \\
\claimstatus{EMP-019}{EMPIRICAL} & \texttt{research/experiments/EXP-0019-m20-iterated-quotient.md} \\
\claimstatus{DEF-021}{DEFINITION} & \texttt{research/proofs/BAR-016-quotient-linear-combination.md} \\
\claimstatus{BAR-016}{PROVED} & \texttt{research/proofs/BAR-016-quotient-linear-combination.md} \\
\claimstatus{REF-017}{REFUTED} & \texttt{research/NEGATIVE\_RESULTS.md NR-018} \\
\claimstatus{EMP-020}{EMPIRICAL} & \texttt{research/experiments/EXP-0020-m21-quotient-linear-combination.md} \\
\claimstatus{DEF-022}{DEFINITION} & \texttt{research/proofs/BAR-017-symmetric-quotient-difference.md} \\
\claimstatus{BAR-017}{PROVED} & \texttt{research/proofs/BAR-017-symmetric-quotient-difference.md} \\
\claimstatus{REF-018}{REFUTED} & \texttt{research/NEGATIVE\_RESULTS.md NR-019} \\
\claimstatus{EMP-021}{EMPIRICAL} & \texttt{research/experiments/EXP-0021-m22-symmetric-quotient-difference.md} \\
\claimstatus{DEF-023}{DEFINITION} & \texttt{research/proofs/BAR-018-unequal-signed-reduction.md} \\
\claimstatus{BAR-018}{PROVED} & \texttt{research/proofs/BAR-018-unequal-signed-reduction.md} \\
\claimstatus{REF-019}{REFUTED} & \texttt{research/NEGATIVE\_RESULTS.md NR-020} \\
\claimstatus{EMP-022}{EMPIRICAL} & \texttt{research/experiments/EXP-0022-m23-unequal-signed-reduction.md} \\
\claimstatus{DEF-024}{DEFINITION} & \texttt{research/proofs/BAR-019-rational-residue-resultants.md} \\
\claimstatus{BAR-019}{PROVED} & \texttt{research/proofs/BAR-019-rational-residue-resultants.md} \\
\claimstatus{REF-020}{REFUTED} & \texttt{research/NEGATIVE\_RESULTS.md NR-021} \\
\claimstatus{EMP-023}{EMPIRICAL} & \texttt{research/experiments/EXP-0023-m24-rational-residue-audit.md} \\
\claimstatus{DEF-025}{DEFINITION} & \texttt{research/proofs/THM-003-rational-root-orbits.md} \\
\claimstatus{THM-003}{PROVED} & \texttt{research/proofs/THM-003-rational-root-orbits.md} \\
\claimstatus{REF-021}{REFUTED} & \texttt{research/NEGATIVE\_RESULTS.md NR-022} \\
\claimstatus{EMP-024}{EMPIRICAL} & \texttt{research/experiments/EXP-0024-m25-rational-root-orbits.md} \\
\claimstatus{DEF-026}{DEFINITION} & \texttt{research/proofs/BAR-020-exceptional-cyclotomic-extraction.md} \\
\claimstatus{BAR-020}{PROVED} & \texttt{research/proofs/BAR-020-exceptional-cyclotomic-extraction.md} \\
\claimstatus{REF-022}{REFUTED} & \texttt{research/NEGATIVE\_RESULTS.md NR-023} \\
\claimstatus{EMP-025}{EMPIRICAL} & \texttt{research/experiments/EXP-0025-m26-exceptional-cyclotomic.md} \\
\claimstatus{DEF-027}{DEFINITION} & \texttt{research/proofs/BAR-021-exceptional-cofactor-schedule-barrier.md} \\
\claimstatus{BAR-021}{PROVED} & \texttt{research/proofs/BAR-021-exceptional-cofactor-schedule-barrier.md} \\
\claimstatus{REF-023}{REFUTED} & \texttt{research/NEGATIVE\_RESULTS.md NR-024} \\
\claimstatus{EMP-026}{EMPIRICAL} & \texttt{research/experiments/EXP-0026-m27-exceptional-cofactor-schedule.md} \\
\claimstatus{DEF-028}{DEFINITION} & \texttt{research/proofs/BAR-022-length-indexed-materialized-support.md} \\
\claimstatus{BAR-022}{PROVED} & \texttt{research/proofs/BAR-022-length-indexed-materialized-support.md} \\
\claimstatus{REF-024}{REFUTED} & \texttt{research/NEGATIVE\_RESULTS.md NR-025} \\
\claimstatus{EMP-027}{EMPIRICAL} & \texttt{research/experiments/EXP-0027-m28-length-indexed-support.md} \\
\claimstatus{DEF-029}{DEFINITION} & \texttt{research/proofs/BAR-023-compact-cofactor-prime-support.md} \\
\claimstatus{BAR-023}{PROVED} & \texttt{research/proofs/BAR-023-compact-cofactor-prime-support.md} \\
\claimstatus{REF-025}{REFUTED} & \texttt{research/NEGATIVE\_RESULTS.md NR-026} \\
\claimstatus{EMP-028}{EMPIRICAL} & \texttt{research/experiments/EXP-0028-m29-compact-cofactor-prime-support.md} \\
\claimstatus{DEF-030}{DEFINITION} & \texttt{research/proofs/BAR-024-compact-support-signatures.md} \\
\claimstatus{BAR-024}{PROVED} & \texttt{research/proofs/BAR-024-compact-support-signatures.md} \\
\claimstatus{REF-026}{REFUTED} & \texttt{research/NEGATIVE\_RESULTS.md NR-027} \\
\claimstatus{EMP-029}{EMPIRICAL} & \texttt{research/experiments/EXP-0029-m30-compact-support-signatures.md} \\
\claimstatus{DEF-031}{DEFINITION} & \texttt{research/proofs/THM-004-BAR-025-diversified-selector.md} \\
\claimstatus{THM-004}{PROVED} & \texttt{research/proofs/THM-004-BAR-025-diversified-selector.md} \\
\claimstatus{BAR-025}{PROVED} & \texttt{research/proofs/THM-004-BAR-025-diversified-selector.md} \\
\claimstatus{REF-027}{REFUTED} & \texttt{research/NEGATIVE\_RESULTS.md NR-028} \\
\claimstatus{EMP-030}{EMPIRICAL} & \texttt{research/experiments/EXP-0030-m31-diversified-compact-signatures.md} \\
\claimstatus{DEF-032}{DEFINITION} & \texttt{research/proofs/THM-005-BAR-026-widened-selector-cap.md} \\
\claimstatus{THM-005}{PROVED} & \texttt{research/proofs/THM-005-BAR-026-widened-selector-cap.md} \\
\claimstatus{BAR-026}{PROVED} & \texttt{research/proofs/THM-005-BAR-026-widened-selector-cap.md} \\
\claimstatus{REF-028}{REFUTED} & \texttt{research/NEGATIVE\_RESULTS.md NR-029} \\
\claimstatus{EMP-031}{EMPIRICAL} & \texttt{research/experiments/EXP-0031-m32-widened-selector-cap.md} \\
\claimstatus{DEF-033}{DEFINITION} & \texttt{research/proofs/THM-006-BAR-027-linear-cap-recurrence.md} \\
\claimstatus{THM-006}{PROVED} & \texttt{research/proofs/THM-006-BAR-027-linear-cap-recurrence.md} \\
\claimstatus{BAR-027}{PROVED} & \texttt{research/proofs/THM-006-BAR-027-linear-cap-recurrence.md} \\
\claimstatus{REF-029}{REFUTED} & \texttt{research/NEGATIVE\_RESULTS.md NR-030} \\
\claimstatus{EMP-032}{EMPIRICAL} & \texttt{research/experiments/EXP-0032-m33-linear-cap-recurrence.md} \\
\claimstatus{THM-007}{PROVED} & \texttt{research/proofs/THM-007-BAR-028-next-envelope.md} \\
\claimstatus{BAR-028}{PROVED} & \texttt{research/proofs/THM-007-BAR-028-next-envelope.md} \\
\claimstatus{REF-030}{REFUTED} & \texttt{research/NEGATIVE\_RESULTS.md NR-031} \\
\claimstatus{EMP-033}{EMPIRICAL} & \texttt{research/experiments/EXP-0033-m34-next-envelope.md} \\
\claimstatus{THM-008}{PROVED} & \texttt{research/proofs/THM-008-BAR-029-length-23-envelope.md} \\
\claimstatus{BAR-029}{PROVED} & \texttt{research/proofs/THM-008-BAR-029-length-23-envelope.md} \\
\claimstatus{REF-031}{REFUTED} & \texttt{research/NEGATIVE\_RESULTS.md NR-032} \\
\claimstatus{EMP-034}{EMPIRICAL} & \texttt{research/experiments/EXP-0034-m35-next-envelope.md} \\
\claimstatus{THM-009}{PROVED} & \texttt{research/proofs/THM-009-BAR-030-length-24-distinct-caps.md} \\
\claimstatus{BAR-030}{PROVED} & \texttt{research/proofs/THM-009-BAR-030-length-24-distinct-caps.md} \\
\claimstatus{REF-032}{REFUTED} & \texttt{research/NEGATIVE\_RESULTS.md NR-033} \\
\claimstatus{EMP-035}{EMPIRICAL} & \texttt{research/experiments/EXP-0035-m36-distinct-cap.md} \\
\claimstatus{THM-010}{PROVED} & \texttt{research/proofs/THM-010-BAR-031-length-25-envelope.md} \\
\claimstatus{BAR-031}{PROVED} & \texttt{research/proofs/THM-010-BAR-031-length-25-envelope.md} \\
\claimstatus{REF-033}{REFUTED} & \texttt{research/NEGATIVE\_RESULTS.md NR-034} \\
\claimstatus{EMP-036}{EMPIRICAL} & \texttt{research/experiments/EXP-0036-m37-length-25-cap.md} \\
\claimstatus{THM-011}{PROVED} & \texttt{research/proofs/THM-011-BAR-032-length-26-envelope.md} \\
\claimstatus{BAR-032}{PROVED} & \texttt{research/proofs/THM-011-BAR-032-length-26-envelope.md} \\
\claimstatus{REF-034}{REFUTED} & \texttt{research/NEGATIVE\_RESULTS.md NR-035} \\
\claimstatus{EMP-037}{EMPIRICAL} & \texttt{research/experiments/EXP-0037-m38-length-26-cap.md} \\
\claimstatus{THM-012}{PROVED} & \texttt{research/proofs/THM-012-BAR-033-length-27-envelope.md} \\
\claimstatus{BAR-033}{PROVED} & \texttt{research/proofs/THM-012-BAR-033-length-27-envelope.md} \\
\claimstatus{REF-035}{REFUTED} & \texttt{research/NEGATIVE\_RESULTS.md NR-036} \\
\claimstatus{EMP-038}{EMPIRICAL} & \texttt{research/experiments/EXP-0038-m39-length-27-cap.md} \\
\claimstatus{THM-013}{PROVED} & \texttt{research/proofs/THM-013-BAR-034-length-28-envelope.md} \\
\claimstatus{BAR-034}{PROVED} & \texttt{research/proofs/THM-013-BAR-034-length-28-envelope.md} \\
\claimstatus{REF-036}{REFUTED} & \texttt{research/NEGATIVE\_RESULTS.md NR-037} \\
\claimstatus{EMP-039}{EMPIRICAL} & \texttt{research/experiments/EXP-0039-m40-length-28-cap.md} \\
\claimstatus{THM-014}{PROVED} & \texttt{research/proofs/THM-014-BAR-035-length-29-envelope.md} \\
\claimstatus{BAR-035}{PROVED} & \texttt{research/proofs/THM-014-BAR-035-length-29-envelope.md} \\
\claimstatus{REF-037}{REFUTED} & \texttt{research/NEGATIVE\_RESULTS.md NR-038} \\
\claimstatus{EMP-040}{EMPIRICAL} & \texttt{research/experiments/EXP-0040-m41-length-29-cap.md} \\
\claimstatus{THM-015}{PROVED} & \texttt{research/proofs/THM-015-BAR-036-length-30-envelope.md} \\
\claimstatus{BAR-036}{PROVED} & \texttt{research/proofs/THM-015-BAR-036-length-30-envelope.md} \\
\claimstatus{REF-038}{REFUTED} & \texttt{research/NEGATIVE\_RESULTS.md NR-039} \\
\claimstatus{EMP-041}{EMPIRICAL} & \texttt{research/experiments/EXP-0041-m42-length-30-cap.md} \\
\claimstatus{THM-016}{PROVED} & \texttt{research/proofs/THM-016-BAR-037-length-31-envelope.md} \\
\claimstatus{BAR-037}{PROVED} & \texttt{research/proofs/THM-016-BAR-037-length-31-envelope.md} \\
\claimstatus{REF-039}{REFUTED} & \texttt{research/NEGATIVE\_RESULTS.md NR-040} \\
\claimstatus{EMP-042}{EMPIRICAL} & \texttt{research/experiments/EXP-0042-m43-length-31-cap.md} \\
\claimstatus{THM-017}{PROVED} & \texttt{research/proofs/THM-017-BAR-038-length-32-envelope.md} \\
\claimstatus{BAR-038}{PROVED} & \texttt{research/proofs/THM-017-BAR-038-length-32-envelope.md} \\
\claimstatus{REF-040}{REFUTED} & \texttt{research/NEGATIVE\_RESULTS.md NR-041} \\
\claimstatus{EMP-043}{EMPIRICAL} & \texttt{research/experiments/EXP-0043-m44-length-32-cap.md} \\
\claimstatus{THM-018}{PROVED} & \texttt{research/proofs/THM-018-BAR-039-length-33-envelope.md} \\
\claimstatus{BAR-039}{PROVED} & \texttt{research/proofs/THM-018-BAR-039-length-33-envelope.md} \\
\claimstatus{REF-041}{REFUTED} & \texttt{research/NEGATIVE\_RESULTS.md NR-042} \\
\claimstatus{EMP-044}{EMPIRICAL} & \texttt{research/experiments/EXP-0044-m45-length-33-cap.md} \\
\claimstatus{THM-019}{PROVED} & \texttt{research/proofs/THM-019-BAR-040-length-34-envelope.md} \\
\claimstatus{BAR-040}{PROVED} & \texttt{research/proofs/THM-019-BAR-040-length-34-envelope.md} \\
\claimstatus{REF-042}{REFUTED} & \texttt{research/NEGATIVE\_RESULTS.md NR-043} \\
\claimstatus{EMP-045}{EMPIRICAL} & \texttt{research/experiments/EXP-0045-m46-length-34-cap.md} \\
\claimstatus{DEF-034}{DEFINITION} & \texttt{research/proofs/BAR-041-polynomial-cap-support-budget.md} \\
\claimstatus{BAR-041}{PROVED} & \texttt{research/proofs/BAR-041-polynomial-cap-support-budget.md} \\
\claimstatus{REF-043}{REFUTED} & \texttt{research/NEGATIVE\_RESULTS.md NR-044} \\
\claimstatus{EMP-046}{EMPIRICAL} & \texttt{research/experiments/EXP-0046-m47-polynomial-cap-support.md} \\
\claimstatus{DEF-035}{DEFINITION} & \texttt{research/proofs/BAR-042-compact-gap-overlap-budget.md} \\
\claimstatus{BAR-042}{PROVED} & \texttt{research/proofs/BAR-042-compact-gap-overlap-budget.md} \\
\claimstatus{REF-044}{REFUTED} & \texttt{research/NEGATIVE\_RESULTS.md NR-045} \\
\claimstatus{EMP-047}{EMPIRICAL} & \texttt{research/experiments/EXP-0047-m48-compact-gap-overlap.md} \\
\claimstatus{DEF-036}{DEFINITION} & \texttt{research/proofs/BAR-043-wide-span-high-weight-overlap.md} \\
\claimstatus{BAR-043}{PROVED} & \texttt{research/proofs/BAR-043-wide-span-high-weight-overlap.md} \\
\claimstatus{REF-045}{REFUTED} & \texttt{research/NEGATIVE\_RESULTS.md NR-046} \\
\claimstatus{EMP-048}{EMPIRICAL} & \texttt{research/experiments/EXP-0048-m49-wide-span-compact-gap.md} \\
\claimstatus{DEF-037}{DEFINITION} & \texttt{research/proofs/BAR-044-subquadratic-span-overlap.md} \\
\claimstatus{BAR-044}{PROVED} & \texttt{research/proofs/BAR-044-subquadratic-span-overlap.md} \\
\claimstatus{REF-046}{REFUTED} & \texttt{research/NEGATIVE\_RESULTS.md NR-047} \\
\claimstatus{EMP-049}{EMPIRICAL} & \texttt{research/experiments/EXP-0049-m51-subquadratic-span.md} \\
\claimstatus{DEF-038}{DEFINITION} & \texttt{research/proofs/BAR-045-boundary-constant-entropy.md} \\
\claimstatus{BAR-045}{PROVED} & \texttt{research/proofs/BAR-045-boundary-constant-entropy.md} \\
\claimstatus{REF-047}{REFUTED} & \texttt{research/NEGATIVE\_RESULTS.md NR-048} \\
\claimstatus{OPEN-004}{REFUTED} & \texttt{research/proofs/BAR-046-distinct-gap-union.md} \\
\claimstatus{EMP-050}{EMPIRICAL} & \texttt{research/experiments/EXP-0050-m52-boundary-constant.md} \\
\claimstatus{DEF-039}{DEFINITION} & \texttt{research/proofs/BAR-046-distinct-gap-union.md} \\
\claimstatus{BAR-046}{PROVED} & \texttt{research/proofs/BAR-046-distinct-gap-union.md} \\
\claimstatus{REF-048}{REFUTED} & \texttt{research/NEGATIVE\_RESULTS.md NR-049} \\
\claimstatus{OPEN-005}{OPEN} & \texttt{research/proofs/BAR-046-distinct-gap-union.md} \\
\claimstatus{EMP-051}{EMPIRICAL} & \texttt{research/experiments/EXP-0051-m53-distinct-gap.md} \\
\bottomrule
\end{longtable}


\end{document}
